\talkslide{objectives}{Learning the defender: online DPO}{Timing: provisional}{Use the defender's own criterion to compare two responses to the same faithful attack.}{
\preferenceflow{1}

}
\speak{This is the same preference-learning construction as for the attacker, now applied to responses. Take one faithful attack and sample two independent responses from the same defender. The defender judge prefers compliance for a benign seed and constructive deflection for a harmful seed. A preference for response $y_2$ supplies a preferred/rejected response pair for online DPO. The judge sees the seed and its class; the defender only sees the attack prompt. Source: AdvGame, Sections 3.1--3.2 and Eq. (8). Without this gate, an attacker could change a benign question into a harmful one and claim success when the defender refuses. The faithfulness judge prevents that shortcut. If only one proposed attack is faithful, it wins. If neither is faithful, there is no useful preference. For the defender, the formal preference is one-half on an unfaithful attack; the practical algorithm skips those defender rollouts. The judge is therefore a real dependency of the method, not a proof that meaning is perfectly preserved.}
